\section{Mordell through splittings}
\label{sec:D}
Let $C/K$ be smooth, proper, geometrically connected of
genus $g\geq2$, with $K$ a number field. Put
$\Delta=\pi_1(C_{\overline K})$, $E=\pi_1(C)$,
$J=\operatorname{Pic}^0_{C/K}$, and
\[
 \operatorname{Sec}(C/K)=
 \operatorname{Sect}_{\rm cts}(E\longrightarrow G_K)/\Delta .
\]
All statements below using the local section theorem retain
the provenance specified in Part B: Theorem 1.1 of the
October 2026 paper is a stipulated input of this problem.
They are deductions from that input and the displayed
published results; they do not claim an independent
verification of every external input of that paper.

\subsection{D(a): exact logical maps and the abelian computation}

\begin{proposition}[D.1 --- Proved: the finiteness diagram]
\label{D-logical}
Grothendieck's assertion on \S8, folio 33, scan 65 is
the finiteness of $\operatorname{Sec}(C/K)$.
The actual maps relevant to it are
\[
 C(K)\ \xrightarrow{\ \kappa\ }\ \operatorname{Sec}(C/K)
 \ \xrightarrow{\ \operatorname{loc}\ }\
 C(\mathbb A_K)_\bullet^{\mathrm{f\text{-}cov}}
 \ \lhook\joinrel\longrightarrow\
 \prod_{v<\infty}C(K_v)\times
 \prod_{v\mid\mathbb R}\pi_0(C(K_v)).
\]
Here $\kappa$ is injective, and $\operatorname{loc}$
is bijective. Consequently the following are exact
equivalences:
\[
 \begin{split}
 \kappa\text{ is bijective}
 &\ \Longleftrightarrow\
 C(\mathbb A_K)_\bullet^{\mathrm{f\text{-}cov}}=C(K),\\
 \#\operatorname{Sec}(C/K)<\infty
 &\ \Longleftrightarrow\
 \#C(\mathbb A_K)_\bullet^{\mathrm{f\text{-}cov}}<\infty.
 \end{split}
\]
The first equality and Faltings's theorem give the second
finiteness. Finiteness of $\operatorname{Sec}(C/K)$ gives
finiteness of $C(K)$ through the displayed injection.
For a finite section set, its cardinality equals that
of $C(K)$ exactly when $\kappa$ is bijective.
\end{proposition}
\begin{proof}
Part B proves Kummer injectivity. The local theorem gives
a unique point at every finite place; Part C gives a
component at every real place. Harari--Stix, Theorem 2.1,
with the normal subgroup $U=1$ and omitted set of places
empty, identifies survival of every finite torsor with
the existence of a global section inducing these local
classes~\cite{HarariStix2012}. Finite torsor evaluations
at real places factor through components. This gives
both inclusions in the asserted image of $\operatorname{loc}$,
exactly as in Proposition 8.1 of the supplied local
paper~\cite{B_OpenAIPadic}.

For injectivity of $\operatorname{loc}$, two sections
with equal local points have equal reductions at all
finite places of good reduction. They are Selmer
sections and are noncuspidal because $C$ is proper.
Porowski's 2025 Theorem 1.1 therefore identifies their
classes~\cite{Porowski2025}; the reduction used in its
application is recorded in Proposition~\ref{D-rigidity}.
The first equivalence can also avoid this last
injectivity theorem: an adelic tuple supported on a
single rational point makes the section a point
section by Stix's finite-support Corollary 5.
The remaining assertions follow from the displayed
injection and bijection. Faltings's theorem says
$\#C(K)<\infty$~\cite{Faltings1983}. We have not reversed
either implication from finiteness to surjectivity;
the actual unresolved part is the possible complement
of the image of $\kappa$, as recorded below.
Stoll's Remark 8.9 establishes the accompanying
section-property, finite-descent, and birational-property
implications in its stated ranges~\cite{Stoll2007}.
\end{proof}


\begin{proposition}[D.2 --- Known, with calculation: every abelian section class]
\label{D-abelian}
The extension with geometric kernel $\Delta^{\rm ab}=TJ$
is a class in $H^2_{\rm cts}(K,TJ)$. Its section set is
empty when that class is nonzero. A choice of one
section identifies its section set with
$H^1_{\rm cts}(K,TJ)$. In particular, a point of $C(K)$
supplies such a choice. There are exact sequences
\[
 0\longrightarrow\widehat{J(K)}
 \longrightarrow H^1_{\rm cts}(K,TJ)
 \longrightarrow T H^1(K,J)\longrightarrow0
\]
and
\[
 0\longrightarrow\widehat{J(K)}
 \longrightarrow\operatorname{Sel}(TJ)
 \longrightarrow T\Sha(K,J)\longrightarrow0.
\]
The Selmer subgroup in the second line is the inverse
image of the local Kummer images at every place.
Here the hats are inverse limits modulo multiplication
by positive integers, and $T$ is the inverse limit
of torsion with the corresponding multiplication
transition maps.
\end{proposition}
\begin{proof}
After choosing a splitting, another splitting has
the form $g\mapsto(c(g),g)$, and the homomorphism
identity is exactly the continuous cocycle identity.
Conjugation by $TJ$ changes $c$ by a coboundary.
Without a chosen splitting the extension cocycle is
the stated degree-two obstruction.

The Kummer exact sequence at level $n$ is
\[
 0\longrightarrow J(K)/nJ(K)
 \longrightarrow H^1(K,J[n])
 \longrightarrow H^1(K,J)[n]\longrightarrow0.
\]
Mordell--Weil makes its left term finite. These terms
have surjective transition maps, so their first
derived inverse limit vanishes. The groups
$H^0(K,J[n])$ are finite as well; thus continuous
cohomology of $TJ$ in degree one is the inverse
limit of these degree-one groups. Taking inverse
limits proves the first line. Intersecting with the
local Kummer conditions at finite level replaces
$H^1(K,J)[n]$ by $\Sha(K,J)[n]$ and gives the second.
This is the Selmer computation used in Stoll,
\S\S3--5~\cite{Stoll2007}.
\end{proof}

\begin{proposition}[D.3 --- Proved: an infinite abelian approximation and an explicit nonpoint class]
\label{D-abelian-examples}
The abelian section set is infinite in general,
even for curves of genus two with a rational base point.
One example is the smooth completion of
\[
 C_+:\quad y^2=x^6-x^2+1.
\]
Moreover, on the genus-two curve
\[
 C_0:\quad y^2=x^5-x+1
\]
with base point $\infty$, the class
\[
 \kappa_J\bigl(2[(0,1)-\infty]\bigr)
 \in H^1(K,TJ),\qquad K=\mathbb Q,
\]
is a global section of the abelianized extension which
is induced by no point of $C_0(K_v)$ at any finite place
$v$.
\end{proposition}
\begin{proof}
The first curve maps to $A:y^2=u^3-u+1$ by
$(x,y)\mapsto(x^2,y)$. The point $P=(0,1)$ on $A$
satisfies $2P=(1/4,-7/8)$. The integral-coordinate
conclusion of the Nagell--Lutz theorem for torsion
on this integral short Weierstrass equation shows
that $P$ has infinite order. The push-pull identity
for the degree-two map is $f_*f^*=[2]$ on $A$.
Consequently $J(C_+)(\mathbb Q)$ has positive rank,
and its completion contains $\widehat{\mathbb Z}$.
Proposition~\ref{D-abelian} proves infinitude.
Both points at infinity of $C_+$ are rational.

For the second curve put $P=(0,1)$. Suppose over
an algebraically closed characteristic-zero field
that $2[P-\infty]=[Q-\infty]$. Then
$2P\sim Q+\infty$ are effective divisors of degree
two. If these divisors differ, they give a moving
degree-two system on a genus-two curve. Riemann--Roch
identifies that system with the unique hyperelliptic
system. But $2P$ is not a hyperelliptic fibre:
$P$ is not a ramification point. If the divisors
are equal, their supports give $P=\infty$, also
impossible. Thus the class is outside the
Abel--Jacobi image over every extension field.

For a finite extension of $\mathbb Q_p$, the group
$J(K_v)$ is compact and has an open subgroup
isomorphic to $\mathbb Z_p^{g[K_v:\mathbb Q_p]}$.
Its power completion is therefore itself, and
its Kummer map into $H^1(K_v,TJ)$ is injective.
A local point inducing the displayed section
would contradict the preceding divisor argument.

The exact program checks smoothness of both
polynomials, the doubling calculation, and the
reduced Mumford representation
\[
 (u,v)=(x^2,1-x/2),\qquad v^2\equiv x^5-x+1\pmod u .
\]
Its reduced degree is two, which independently
certifies the degree-one Abel--Jacobi exclusion.
These are statements about the actual quotient
$\Delta\to TJ$, with its induced section map.
\end{proof}

\subsection{D(b): the finite-descent criterion and explicit curves}

\begin{theorem}[D.4 --- Proved: the Stoll range in terms of splittings]
\label{D-Stoll}
Every section is induced by a unique rational point
for every proper hyperbolic curve $C/K$ with a
nonconstant map to an abelian variety $A/K$ satisfying
\[
 \#A(K)<\infty,\qquad
 \Sha(K,A)_{\rm div}=0 .
\]
In particular this holds for $X_0(N)$ and $X_1(N)$
over $\mathbb Q$ of genus at least two. The section
set is finite in each case. Explicit genus-two
examples over $\mathbb Q$ are the completions of
\[
 y^2=x^6+a,\qquad a\in\{-2,-1,1,2\}.
\]
Explicit genus-three examples are
\[
 y^2=x^8+4,\qquad y^2=x^8+64.
\]
\end{theorem}
\begin{proof}
Stoll's abelian descent computation identifies
$A(\mathbb A_K)_\bullet^{\rm f\text{-}ab}$ with $A(K)$
under the displayed assumptions. Push a finite-descent
adelic point of $C$ to $A$. It lands at a point
$a\in A(K)$. The fibre $Z=f^{-1}(a)$ is finite,
since $f$ is nonconstant and $C$ is a proper curve.
Stoll's finite-support theorem, Theorem 8.2, gives
\[
 C(\mathbb A_K)_\bullet^{\rm f\text{-}ab}
 \cap Z(\mathbb A_K)_\bullet=Z(K).
\]
This proves the finite-descent equality for $C$,
and Proposition~\ref{D-logical} gives precisely
the required splitting statement. We use the
corrected real-component convention in Stoll's
erratum, \S3; finite-place support is sufficient
for this argument~\cite{StollErratum}.

For the modular curves, Corollary 8.8 of Stoll
supplies this equality: the Jacobian has a
nontrivial modular quotient of analytic rank zero,
and the rank-zero finiteness results give the
required Mordell--Weil and Tate--Shafarevich
properties. This is the complete chain behind
Corollary 8.5 of the supplied 2026 local paper:
local Theorem 1.1, Harari--Stix Theorem 2.1,
Stix Corollary 5, and Stoll Corollary 8.8.

For $y^2=x^6+a$, the two maps to elliptic curves are
\[
 (x,y)\longmapsto(x^2,y)\in E_a,\qquad
 (x,y)\longmapsto(a/x^2,ay/x^3)\in E_{a^2},
 \quad E_b:\ v^2=u^3+b .
\]
They extend to the smooth completions. Stoll's
examples immediately following Theorem 8.6
verify the analytic-rank-zero criterion for
the four listed values. For genus three use
$A:y^2=x^3+4x$, the elliptic modular curve
$X_0(32)$, whose rational Mordell--Weil group has
rank zero and whose Tate--Shafarevich group is finite.
The maps are respectively
\[
 (x,y)\longmapsto(x^4,x^2y),\qquad
 (x,y)\longmapsto(x^4/4,x^2y/8).
\]
Direct substitution verifies their equations.
The squarefree polynomials of degree eight
give genus three. The identification of this
elliptic curve with $X_0(32)$, or its standard
degree-two isogeny to $y^2=x^3-x$, supplies
the same rank-zero modular input~\cite{X032Model}.
Finally Faltings's theorem makes each rational
point set finite.
\end{proof}

\subsection{D(c): finite local images and the global fibre}

\begin{proposition}[D.5 --- Known, with checked reduction: local finiteness and reduction rigidity]
\label{D-rigidity}
For $K$ containing no CM subfield, the image
$\operatorname{loc}_v\operatorname{Sec}(C/K)$
is finite at every finite place $v$.
For arbitrary number fields $K$, the reduction
of Selmer sections on a set of strongly positive
density determines their classes. Strongly
positive density means positive upper density
after pullback to the primes of every finite
extension of $K$.
\end{proposition}
\begin{proof}
Every global section is Selmer by the stipulated
local theorem. Betts--Stix, Theorem A, applies
to exactly this Selmer set and gives the first
claim~\cite{BettsStix2025}. Its input is a finite
set of global semisimple Galois representations
from the chosen abelian-by-finite family, followed
by the finite period-map loci of their types
(a) and (b); type (c) is excluded. The proof uses
Corollaries 6.10 and 6.16 and Proposition 6.20
to obtain Theorem 6.5. The choice of auxiliary
family depends on the place. No place-independent
cardinality bound is asserted there.

For the second claim we checked the following
chain in the supplied source of Porowski 2026,
Theorem 1.5~\cite{Porowski2026}. A finite
semistable partition at a fixed place is
separated, after a finite cover and a finite
constant extension, by finitely many smooth
proper families whose Jacobians have different
toric ranks across distinct parts. The
construction is the source's Theorem 6.4,
its compactness step in \S5, and the
resolution-of-nonsingularities step in \S6.
For the actual families the semisimple
representations are integral, pure of weight
one, of bounded dimension, and unramified
outside a fixed finite set. The Faltings
finiteness theorem used in \S7 therefore
gives a finite list.

Pass to a finite extension over which the
connected algebraic monodromy groups of each
pair are connected. If two such semisimple
representations are not potentially isomorphic,
their trace-equality locus is a proper
algebraic subset of that connected monodromy
group, hence has Haar measure zero. This is
the algebraic Chebotarev argument in Rajan's
Theorem 2~\cite{Rajan1998}. Intersecting the
finitely many complementary sets produces a
set of potentially density one. Equality of
reductions at any prime in that set forces
equality of the fixed local partition, since
different toric ranks remain different under
finite local extension. This is precisely
the source's set $\Omega(v)$.

A strongly positive set meets every potentially
density-one set: pull back to an extension
where the latter has density one and use
positive upper density of the former.
Consequently equality of reductions on the
given strongly positive set gives equality
at every place. Porowski 2025, Theorem 1.1,
then identifies the sections. Its density-one
argument uses the weak-cover property of
Theorem 3.3 and the noncuspidal reduction
estimate in Theorem 5.4; the proper curve
here has no boundary and satisfies their
noncuspidal hypotheses. This records the
exact use of both recent papers. We use no
extension of the density hypothesis to
arbitrary infinite sets.
\end{proof}

\begin{proposition}[D.6 --- Proved: the precise uniform reduction criterion]
\label{D-uniform}
Put $r(v)=\#\operatorname{red}_v\operatorname{Sel}(C/K)$.
For each integer $d\geq1$, the inequality
$\#\operatorname{Sel}(C/K)\leq d$ is equivalent
to the existence of a strongly positive set
$\Omega$ with $r(v)\leq d$ for every $v\in\Omega$.
The step in the supplied 2026 source replacing
strongly positive density by strictly positive
density has an unproved intersection assertion.
\end{proposition}
\begin{proof}
The forward implication uses all finite places.
For the reverse, take $d+1$ distinct classes.
At every $v\in\Omega$ some pair has equal
reduction. Thus $\Omega$ is a union of
finitely many pair-equality sets. At least
one of those sets has strongly positive
density. Indeed, for each set failing
strong positivity there is a finite extension
where its pullback has upper density zero.
The compositum of these finitely many
extensions makes all their pullbacks have
upper density zero; finite extension preserves
density zero, since degree-one primes have
bounded multiplicity and primes of residue
degree greater than one have density zero.
Their union cannot contain the pullback
of $\Omega$. Apply Proposition~\ref{D-rigidity}
to the resulting pair to get a contradiction.

The unnumbered corollary following the source's
Corollary 7.10 states the bound with strictly
positive density, and its proof chooses
$v'\in\Omega\cap\Omega(v)$. The established
property of $\Omega(v)$ is potentially density
one. For a nontrivial quadratic extension
$L/K$, the inert primes have density $1/2$,
the completely split primes have potentially
density one, and their intersection is empty.
This explicit Chebotarev example refutes the
intersection inference in that proof.
It does not supply a counterexample curve
to the corollary itself. The strongly positive
version has the complete proof above.
\end{proof}

\begin{proposition}[D.7 --- Proved: an arithmetic model of the remaining uniformity obstruction]
\label{D-density-model}
There is an infinite compact totally disconnected
Hausdorff space with continuous finite images at
each rational prime, with injective product
map over every strongly positive set of primes,
and with all the potentially-density-one
partition refinements used in the preceding
proof. Thus these three properties, as
properties of continuous maps of compact spaces, do not establish
finiteness of the source set.
\end{proposition}
\begin{proof}
Enumerate the rational primes $p_1,p_2,\ldots$,
and put $L_m=\mathbb Q(\sqrt{p_1},\ldots,\sqrt{p_m})$,
$L_0=\mathbb Q$. Let $b(p)$ be the largest
$m$ for which $p$ is unramified and completely
split in $L_m$. It is finite, because the
step adjoining $\sqrt p$ ramifies at $p$.
Give $S=\mathbb N_0\cup\{\infty\}$ the one-point
compactification topology, and define
\[
 f_p(n)=\min\{n,b(p)\},\qquad f_p(\infty)=b(p).
\]
Every image is finite and discrete. Each map is
continuous, because it is constant on the cofinite
neighbourhood $\{n\ge b(p)\}\cup\{\infty\}$.
This proves the required compact topology explicitly.
For any distinct finite $n,n'$
and any strongly positive $\Omega$, the set
of primes completely split in
$L_{\max(n,n')}$ has potentially density one,
so it meets $\Omega$. At such a prime
$f_p(n)\ne f_p(n')$. The product is injective.
For the pair $n,\infty$, use instead the completely
split primes of $L_{n+1}$; there $b(p)>n$, so
$f_p(n)\ne f_p(\infty)$ as well. The continuous
product injection is a homeomorphism onto a closed
subspace of the product of its finite images,
since its source is compact and the target Hausdorff.

For a fixed $v$, complete splitting in
$L_{b(v)}$ implies that equality under $f_p$
implies equality under $f_v$. Hence the
partition-refining set has potentially
density one. All assertions are realized
by the explicit maps $f_p$, while
$S$ is infinite. Its sequence $n\to\infty$
also realizes a compatible limit that every fixed
finite collection of local maps eventually fails
to distinguish. The program
records the first splitting tests using
Legendre symbols. We make no identification
of this set with a section set of a curve.
The morphisms displayed here prove exactly
which purely formal inference fails.
\end{proof}

\begin{proposition}[D.8 --- Open: the global finiteness fibre]
\label{D-open}
The finiteness of $\operatorname{Sec}(C/K)$
for every proper hyperbolic curve over a number
field remains undecided in this return.
Proposition~\ref{D-literal-minimum} gives the literal
least quotient for every locally soluble curve;
Theorem~\ref{D-no-least-quotient} disproves the existence
of a least quotient in general.
\end{proposition}
\begin{proof}[State of the attack]
The maps in Proposition~\ref{D-logical} identify
the first issue with finiteness of the
finite-descent locus. Proposition~\ref{D-rigidity}
makes each localization image finite over
fields without CM subfields, while
Proposition~\ref{D-uniform} identifies a
sufficient and necessary uniform reduction
bound. Proposition~\ref{D-density-model}
proves why the placewise and density
statements alone do not supply it.

We tried four routes. The finite list of
Galois representations in a period-map
argument bounds the local image for its
chosen auxiliary family; the degree of the
separating cover and the family change
when another local partition is imposed.
No bounded common list was obtained.
The finite-cover neighbourhood argument
turns a section supported at a known
rational point into that point section,
but it supplies no such support for a
general member of a localization fibre.
The potentially-density-one refinements
control each finite partition separately;
the explicit model above shows the missing
uniform bound. Finally, compactness of an
inverse system of finite quotient data
constructs compatible limits; it gives
no finite bound on their number.

The smallest genus left by this argument
is two over $\mathbb Q$, after the
finite-descent cases of Theorem~\ref{D-Stoll}
have been removed. More precisely the
unresolved object is a fibre of
\[
 \operatorname{Sec}(C/\mathbb Q)
 \longrightarrow C(\mathbb Q_p)
\]
over one of its finitely many values, for
a genus-two curve to which no verified
finite-descent criterion has been applied.
The next mathematical step is to bound
that fibre using a fixed finite collection
of separating families, or to prove a
uniform bound as in Proposition~\ref{D-uniform};
neither bound is asserted here.

For the quotient question, quotients are
ordered by factorization of continuous
surjections over $G_K$, and local
point-inducedness means membership in
the image of the actual map
$C(K_v)\to\operatorname{Sec}(E_Q/G_{K_v})$.
The full group has the required property
under the stipulated local input.
Proposition~\ref{D-abelian-examples}
excludes the abelian quotient by an
explicit global section. The cited
local pro-prime, metabelian and
nilpotent-by-abelian counterexamples
cannot by themselves produce a global
section: the required map is the
restriction
$\operatorname{Sec}(E_Q/G_K)\to
\operatorname{Sec}(E_Q/G_{K_v})$,
whose image must be controlled.
The independent global construction
below addresses one such image.
Pop's small birational quotients start
with $G_{K(C)}$ and include valuative
inertia; the map to $\pi_1(C)$ kills
the inertia of every point of $C$.
We therefore retain the actual quotient
maps and do not infer a least quotient
of $\pi_1(C)$ from that birational theorem.
\end{proof}


\begin{theorem}[D.9 --- Proved: explicit global pro-three sections failing local geometricity]
\label{D-global-prothree}
Put
\[
 L=\mathbb Q(\zeta_9,\sqrt[3]{3}),\qquad
 X_9:\ X^9+Y^9+Z^9=0 .
\]
The proper curve $X_9/L$ has genus $28$. Its
geometrically pro-$3$ arithmetic fundamental
group has infinitely many section classes
which are induced by no $L_w$-point at a
place $w$ above $3$. The same construction
gives such sections after passing to each
of its geometric pro-$3$ nilpotent or
solvable quotients of positive class or
derived length.

This is an explicit instance of the global
phenomenon in Hoshi's 2010 Theorems A and B,
whose proof and scope we checked; we make
no priority claim~\cite{Hoshi2010}.
\end{theorem}
\begin{proof}
Let $K=\mathbb Q(\zeta_3)$ and let $G$ be
the maximal pro-$3$ Galois group unramified
outside its unique prime above $3$.
The field has class number one, contains
$\mu_3$, and has just one prime above $3$.
Shafarevich's theorem makes $G$ free pro-$3$;
the precise reduction is given in
Hoshi, Lemma 3.2, and in
\c{C}iperiani--Stix, \S4, using
NSW Corollary 10.7.14~\cite{CiperianiStix}.
Its rank is two: Kummer theory identifies
$H^1(G,\mathbb F_3)$ with the two-dimensional
space of $S$-units modulo cubes, generated
by $\zeta_3$ and $\lambda=1-\zeta_3$.

Let $E=X_3$, based at $O=[1:-1:0]$.
Its equation is $y^2+y=x^3-7$, or
$Y^2=X^3-432$. All of $E[3]$ is rational
over $K$: the nonzero division points have
$X=0$, $Y=\pm12\sqrt{-3}$, or
$X=12\zeta_3^i$, $Y=\pm36$.
The curve has good reduction outside $3$.
Thus the action on $M=T_3E$ factors through $G$.

For completeness the rank-zero arithmetic input
has an exact verification. The normalized
newform of $E$ is
\[
 f(z)=\eta(3z)^2\eta(9z)^2 ,
\]
as recorded in Alfes--Mertens, \S4.2~\cite{AlfesMertens2024}.
Every factor in its product expansion is
positive on the positive imaginary axis, so
\[
 L(E,1)=2\pi\int_0^\infty f(it)\,dt>0.
\]
The integral converges at both cusps. The
quadratic twist by $-3$ has equation
$Y^2=X^3+11664$ and is rationally
$3$-isogenous to $E$, through the standard
map from $Y^2=X^3+B$ to
$Y^2=X^3-27B$. It has the same nonzero
$L$-value. The rank-zero theorem of
Kolyvagin gives finite Mordell--Weil and
Tate--Shafarevich groups for both curves
over $\mathbb Q$~\cite{Kolyvagin1988}.
The plus and minus decomposition for
$K/\mathbb Q$, whose degree is prime to
$3$, gives
\[
 \operatorname{rank}E(K)=0,\qquad
 \#\Sha(K,E)[3^\infty]<\infty .
\]

Choose free generators $\tau_1,\tau_2$ of
$G$. Continuous cocycles into $M$ are
arbitrary pairs of elements of $M$,
and their coboundaries are
\[
 m\longmapsto
 \bigl((\tau_1-1)m,(\tau_2-1)m\bigr).
\]
The kernel is $M^G=0$, since $E(K)$ has
finite torsion. Hence $H^1(G,M)$ has
$\mathbb Z_3$-rank two. Its image in
$H^1(K,M)$ is injective by inflation.
The global $3$-adic Selmer subgroup is
finite by the rank and Tate--Shafarevich
calculation. Choose an infinite-order
class $a\in H^1(G,M)$.

At a finite place away from $3$, the
unramified cohomology of the good elliptic
curve is exactly its local Kummer image.
This follows from the good-reduction
Kummer sequence and
$M/(\operatorname{Frob}-1)M$; it is also
the calculation in
\c{C}iperiani--Stix, \S2.2.
There is just one remaining finite place
of $K$. Consequently the image
\[
 \overline a_3\in T_3H^1(K_3,E)
\]
is nonzero. This explicitly locates the
failure of local geometricity.

Now put $U_3=X_3\setminus\{XYZ=0\}$ and
$N=\pi_1(U_{3,\overline K})^{(3)}$.
The group $N$ is free pro-$3$ of rank ten.
Use the tangential section at $O$ with
integral parameter $t=Z/X$ and tangent
value one. The based action on $N$ is
unramified outside $3$: the smooth pair
and this nonzero tangent extend over
$\mathcal O_K[1/3]$, and the prime-to-residue-
characteristic specialization equivalence
preserves the chosen punctured-disc fibre
functor. This establishes the assertion
for the actual action.

On $N^{\rm ab}/3$, its submodule of
boundary classes and its elliptic
quotient both have trivial $G_K$-action:
the nine cusps and $E[3]$ are $K$-rational,
and $\mu_3\subset K$. Its residual image
is therefore a group of unipotent block
matrices of order a power of three.
The kernel of
$\operatorname{Aut}(N)\to
\operatorname{Aut}(N^{\rm ab}/3)$
is pro-$3$, by the characteristic
Frattini series. The full based action
has pro-$3$ image and factors through $G$.

The graph of a representative of $a$ is
a homomorphism $G\to M\rtimes G$.
Projectivity of the free pro-$3$ group
lifts it through the epimorphism
\[
 N\rtimes G\longrightarrow M\rtimes G.
\]
The projection of the lift to $G$ is
the identity, because the displayed
map preserves it. Pulling back along
$G_K\to G$ gives a section of the
geometrically pro-$3$ group of $U_3$,
whose image on $E$ is $a$.

The coordinate-cube map $U_9\to U_3$
is a degree-nine torsor under $\mu_3^2$.
It is defined over $K$. At the chosen
cusp its parameters satisfy $t_3=t_9^3$,
so the reference tangential section has
zero covering character. The covering
character of the newly constructed
section factors through
$G^{\rm ab}/3$. Its fixed extension is
contained in
\[
 K(\sqrt[3]{\zeta_3},\sqrt[3]{\lambda})
   =K(\zeta_9,\sqrt[3]{3})=L ,
\]
where $\lambda^2=-3\zeta_3$ verifies
the equality. Thus over this fixed
field $L$ the section lifts to $U_9$.
The geometric subgroup corresponding
to this cover is an open subgroup of
$N$, and pro-$3$ covering comparison
identifies it with the pro-$3$ group
of $U_9$. Filling its cusps gives a
section on $X_9$.

Restriction on $T_3H^1(K_3,E)$ is
injective after any finite local
extension: corestriction composed
with restriction is multiplication
by the extension degree, and this
Tate module is $\mathbb Z_3$-torsion-free.
The image of $\overline a_3$ in
$T_3H^1(L_w,E)$ remains nonzero.
A point inducing the constructed
section on $X_9/L_w$ would map under
$X_9\to X_3=E$ to a point inducing
that elliptic section, a contradiction.

The rank-two lattice modulo its
finite Selmer subgroup provides
infinitely many distinct local
non-Kummer images, and restriction
preserves them. Hence infinitely
many classes remain distinct on
$X_9$. Finally every positive
nilpotent-class or derived-length
quotient still maps to the abelian
group $T_3E$, so the same obstruction
survives under the stated quotient
maps. The genus is
$(9-1)(9-2)/2=28$.
The program checks the elliptic
equations, division polynomial,
isogeny identity, and covering ranks.
\end{proof}


